\documentclass[a4paper, 12pt]{article}
\usepackage[left=3cm, top=3cm, right=2cm, bottom=2cm]{geometry}
\usepackage[utf8]{inputenc}
\usepackage[T1]{fontenc}
\usepackage[brazil]{babel}
\usepackage{graphicx}
\usepackage{fvextra}
\pagestyle{myheadings}
\usepackage{setspace}
\usepackage{indentfirst}
\setlength{\parindent}{1cm}
\usepackage{float}
\usepackage{amsthm}
\usepackage{amsmath}
\usepackage{amsfonts}
\usepackage{amssymb}
\usepackage{rotating}
\usepackage{caption}
\usepackage{url}
\usepackage{hyperref}
\usepackage{listings}
\usepackage{xcolor}
\usepackage{upquote}

\lstset{
    basicstyle=\ttfamily\small,
    frame=single,
    breaklines=true,
    tabsize=4,
    showstringspaces=false,
    columns=fullflexible,
    keepspaces=true
}

\lstdefinestyle{minijava}{
    keywordstyle=\color{blue}\bfseries,
    commentstyle=\color{gray}\itshape,
    stringstyle=\color{red},
    numbers=left,
    numberstyle=\tiny\color{gray},
    morekeywords={class, public, static, void, main, String, return, int, boolean, if, else, while, System, out, println, true, false, this, new, extends, length}
}

\lstdefinestyle{cpp}{
    language=C++,
    keywordstyle=\color{blue}\bfseries,
    commentstyle=\color{gray}\itshape,
    stringstyle=\color{red},
    numbers=left,
    numberstyle=\tiny\color{gray}
}

\begin{document}

%CAPA

\begin{center}
    
%FOLHA DE ROSTO
    \setcounter{page}{1} %comando que redefine a contagem de páginas
    \thispagestyle{empty}

    JOÃO GUILHERME LOPES ALVES DA COSTA\\
    JOSÉ DAVI VIANA FRANCELINO\\
    JUSCELINO PEREIRA DE ARAUJO\\
    THIAGO DE OLIVEIRA CORDEIRO\\
    VICTOR BASTOS XAVIER\\
    \vspace{5.2cm}
    
    \textbf{COMPILADOR MINIJAVA: IMPLEMENTAÇÃO DE UM COMPILADOR COM ANÁLISE LÉXICA, SINTÁTICA E TABELA DE SÍMBOLOS}\\
    \vspace{5.2cm}
    
\end{center}

\hspace{8cm} % posicionando a caixa de texto
\begin{minipage}{7cm}
Relatório técnico apresentado como avaliação 1 da disciplina DIM0164 — Compiladores, ministrada pelo(a) professor(a) Valdigleis da Silva Costa, para o curso de Ciência da Computação do Departamento de Informática e Matemática Aplicada da Universidade Federal do Rio Grande do Norte - Campus Natal. 
\end{minipage}
\vspace{4.5cm}

\begin{center}
    NATAL\\
    2026
\end{center}

\singlespacing

\newpage
\tableofcontents
\thispagestyle{empty}

\section{Introdu\c{c}\~ao}

O presente relat\'orio tem como objetivo apresentar a documenta\c{c}\~ao do compilador para a linguagem \textit{MiniJava}, desenvolvido como projeto da disciplina DIM0164 --- Compiladores (2026.1, UFRN). O compilador implementa as fases de an\'alise l\'exica, an\'alise sint\'atica e constru\c{c}\~ao da tabela de s\'imbolos.

A linguagem \textit{MiniJava} \'e um subconjunto da linguagem Java, projetada para fins did\'aticos. O compilador foi implementado em C++ (padr\~ao C++17) utilizando o gerador de analisadores l\'exicos \textit{Flex} e um parser recursivo descendente escrito manualmente para valida\c{c}\~ao da gram\'atica LL(1).

\section{A linguagem MiniJava}

\subsection{Vis\~ao geral}

MiniJava \'e uma linguagem orientada a objetos com tipagem est\'atica. Todo programa MiniJava possui obrigatoriamente uma classe principal com um m\'etodo \texttt{main}, seguida de zero ou mais classes auxiliares.

\subsection{Tipos de dados}

A linguagem suporta os seguintes tipos:

\begin{itemize}
    \item \texttt{int} --- inteiros;
    \item \texttt{int[]} --- arrays de inteiros;
    \item \texttt{boolean} --- valores l\'ogicos (\texttt{true} ou \texttt{false});
    \item Tipos definidos pelo usu\'ario (nomes de classes).
\end{itemize}

\subsection{Estrutura de um programa}

Todo programa MiniJava segue esta estrutura:

\begin{lstlisting}[style=minijava, caption={Estrutura b\'asica de um programa MiniJava}]
class NomePrincipal {
    public static void main(String[] args) {
        // comandos do programa principal
    }
}

class OutraClasse {
    // variaveis de instancia
    int x;

    // metodos
    public int metodo(int param) {
        // corpo
        return valor;
    }
}
\end{lstlisting}

\subsection{Declara\c{c}\~ao de vari\'aveis}

Vari\'aveis s\~ao declaradas com tipo seguido de nome e ponto-e-v\'irgula:

\begin{lstlisting}[style=minijava, caption={Declara\c{c}\~ao de vari\'aveis}]
int x;
int[] numeros;
boolean flag;
MinhaClasse obj;
\end{lstlisting}

\subsection{Comandos}

\subsubsection{Atribui\c{c}\~ao}

\begin{lstlisting}[style=minijava, caption={Atribui\c{c}\~oes}]
x = 10;
numeros[0] = 42;
\end{lstlisting}

\subsubsection{Condicional (if/else)}

O \texttt{if} em MiniJava sempre exige um bloco \texttt{else}:

\begin{lstlisting}[style=minijava, caption={Estrutura condicional}]
if (x < 10)
    y = 1;
else
    y = 2;
\end{lstlisting}

\subsubsection{La\c{c}o de repeti\c{c}\~ao (while)}

\begin{lstlisting}[style=minijava, caption={La\c{c}o while}]
while (i < n) {
    System.out.println(i);
    i = i + 1;
}
\end{lstlisting}

\subsubsection{Impress\~ao}

A \'unica forma de sa\'ida \'e \texttt{System.out.println}, que imprime um inteiro:

\begin{lstlisting}[style=minijava, caption={Impress\~ao de valores}]
System.out.println(42);
System.out.println(x + y);
\end{lstlisting}

\subsection{Express\~oes e operadores}

\subsubsection{Operadores aritm\'eticos}

\begin{itemize}
    \item \texttt{+} (adi\c{c}\~ao), \texttt{-} (subtra\c{c}\~ao), \texttt{*} (multiplica\c{c}\~ao)
\end{itemize}

\subsubsection{Operadores de compara\c{c}\~ao e l\'ogicos}

\begin{itemize}
    \item \texttt{<} (menor que), \texttt{>} (maior que)
    \item \texttt{\&\&} (E l\'ogico), \texttt{!} (nega\c{c}\~ao)
\end{itemize}

\subsubsection{Acesso a arrays e m\'etodos}

\begin{lstlisting}[style=minijava, caption={Express\~oes com arrays e m\'etodos}]
// Acesso a array
x = numeros[i];

// Tamanho do array
tam = numeros.length;

// Chamada de metodo
resultado = obj.metodo(arg1, arg2);
\end{lstlisting}

\subsection{Classes e heran\c{c}a}

Classes podem herdar de outra classe usando \texttt{extends}:

\begin{lstlisting}[style=minijava, caption={Heran\c{c}a de classes}]
class Animal {
    int idade;

    public int getIdade() {
        return idade;
    }
}

class Cachorro extends Animal {
    int peso;

    public int getPeso() {
        return peso;
    }
}
\end{lstlisting}

\subsection{Instancia\c{c}\~ao de objetos e arrays}

\begin{lstlisting}[style=minijava, caption={Cria\c{c}\~ao de objetos e arrays}]
// Novo objeto
obj = new MinhaClasse();

// Novo array de inteiros com 10 posicoes
numeros = new int[10];
\end{lstlisting}

\subsection{Exemplo completo: Fatorial}

\begin{lstlisting}[style=minijava, caption={Programa que calcula o fatorial de 10}]
class Factorial {
    public static void main(String[] a) {
        System.out.println(new Fac().ComputeFac(10));
    }
}

class Fac {
    public int ComputeFac(int num) {
        int num_aux;
        if (num < 1)
            num_aux = 1;
        else
            num_aux = num * (this.ComputeFac(num - 1));
        return num_aux;
    }
}
\end{lstlisting}

Neste exemplo:
\begin{itemize}
    \item A classe \texttt{Factorial} cont\'em o \texttt{main}, que instancia \texttt{Fac} e chama \texttt{ComputeFac};
    \item A classe \texttt{Fac} define um m\'etodo recursivo que calcula o fatorial;
    \item \texttt{this.ComputeFac(num - 1)} \'e uma chamada recursiva ao pr\'oprio m\'etodo.
\end{itemize}

\subsection{Coment\'arios}

A linguagem suporta dois estilos de coment\'arios:

\begin{lstlisting}[style=minijava, caption={Coment\'arios}]
// Comentario de linha unica

/* Comentario de
   multiplas linhas */
\end{lstlisting}

\section{\textit{Design} da implementa\c{c}\~ao}

\subsection{Etapas de implementa\c{c}\~ao}

O desenvolvimento do compilador seguiu as seguintes etapas:

\begin{enumerate}
    \item \textbf{Transforma\c{c}\~ao da gram\'atica para LL(1)}: a gram\'atica original da linguagem MiniJava foi transformada para eliminar recurs\~ao \`a esquerda e fatorar produ\c{c}\~oes amb\'iguas, resultando em uma gram\'atica LL(1) adequada para um parser recursivo descendente;
    \item \textbf{Implementa\c{c}\~ao do analisador l\'exico}: defini\c{c}\~ao dos tokens e constru\c{c}\~ao do scanner com \textit{Flex};
    \item \textbf{Implementa\c{c}\~ao do analisador sint\'atico}: constru\c{c}\~ao do parser recursivo descendente a partir da gram\'atica LL(1);
    \item \textbf{Constru\c{c}\~ao da tabela de s\'imbolos}: integra\c{c}\~ao ao parser para registrar declara\c{c}\~oes durante o parsing.
\end{enumerate}

\subsection{Transforma\c{c}\~ao para gram\'atica LL(1)}

A gram\'atica original da linguagem MiniJava possui constru\c{c}\~oes que n\~ao s\~ao diretamente compat\'iveis com parsers LL(1), como produ\c{c}\~oes com prefixos comuns. Para resolver isso, foram aplicadas as seguintes t\'ecnicas:

\begin{itemize}
    \item \textbf{Fatora\c{c}\~ao \`a esquerda}: produ\c{c}\~oes como \texttt{Cmd}, que podem come\c{c}ar com \texttt{Id} tanto para atribui\c{c}\~ao simples quanto para atribui\c{c}\~ao em array, foram fatoradas em \texttt{Cmd} e \texttt{Cmd'};
    \item \textbf{Elimina\c{c}\~ao de recurs\~ao \`a esquerda}: express\~oes que tinham recurs\~ao \`a esquerda (como opera\c{c}\~oes bin\'arias encadeadas) foram transformadas usando o padr\~ao \texttt{Exp} $\rightarrow$ \texttt{PrimExp Exp'} com \texttt{Exp'} recursivo \`a direita;
    \item \textbf{Separa\c{c}\~ao de alternativas}: produ\c{c}\~oes como \texttt{DefCl'} separam o caso com heran\c{c}a (\texttt{extends}) do caso sem.
\end{itemize}

A gram\'atica LL(1) resultante \'e apresentada a seguir:

\begin{lstlisting}[numbers=none, escapeinside={(*}{*)}]
Prog      (*$\rightarrow$*) MainC DefCl

MainC     (*$\rightarrow$*) 'class' Id '{' 'public' 'static' 'void' 'main'
            '(' 'String' '[' ']' Id ')' '{' Cmd '}' '}'

DefCl     (*$\rightarrow$*) 'class' Id DefCl' | (*$\lambda$*)
DefCl'    (*$\rightarrow$*) '{' DefVar DefMet '}' DefCl
          | 'extends' Id '{' DefVar DefMet '}' DefCl

DefVar    (*$\rightarrow$*) Type Id ';' DefVar | (*$\lambda$*)

DefMet    (*$\rightarrow$*) 'public' Type Id '(' DefMetArgs ')'
            '{' DefVar Cmd 'return' Exp ';' '}' DefMet | (*$\lambda$*)
DefMetArgs (*$\rightarrow$*) Args | (*$\lambda$*)

Type      (*$\rightarrow$*) 'int' '[' ']' | 'boolean' | 'int' | Id

Args      (*$\rightarrow$*) Type Id Args'
Args'     (*$\rightarrow$*) ',' Args | (*$\lambda$*)

Cmd       (*$\rightarrow$*) '{' CmdList '}'
          | 'if' '(' Exp ')' Cmd 'else' Cmd
          | 'while' '(' Exp ')' Cmd
          | 'System' '.' 'out' '.' 'println' '(' Exp ')' ';'
          | Id Cmd'
Cmd'      (*$\rightarrow$*) '=' Exp ';'
          | '[' Exp ']' '=' Exp ';'
CmdList   (*$\rightarrow$*) Cmd CmdList | (*$\lambda$*)

Exp       (*$\rightarrow$*) PrimExp Exp'
Exp'      (*$\rightarrow$*) Op PrimExp Exp'
          | '[' Exp ']' Exp'
          | '.' PostExpDot
          | (*$\lambda$*)

Op        (*$\rightarrow$*) '&&' | '<' | '>' | '+' | '-' | '*'

PostExpDot (*$\rightarrow$*) 'length' Exp'
           | Id '(' ListExp ')' Exp'

ListExp   (*$\rightarrow$*) Exp ListExp' | (*$\lambda$*)
ListExp'  (*$\rightarrow$*) ',' Exp ListExp' | (*$\lambda$*)

PrimExp   (*$\rightarrow$*) 'new' PrimExp' | '!' Exp | '(' Exp ')'
          | 'true' | 'false' | Id | Number | 'this'
PrimExp'  (*$\rightarrow$*) Id '(' ')' | 'int' '[' Exp ']'
\end{lstlisting}

Esta gram\'atica garante que, para cada n\~ao-terminal, o token atual (lookahead) \'e suficiente para determinar univocamente qual produ\c{c}\~ao aplicar.

\subsection{Estrutura do projeto}

\begin{lstlisting}[numbers=none]
src/
  token.h          - Definicao dos tipos de token
  lexer.l          - Analisador lexico (Flex)
  parser.h         - Analisador sintatico
  symbol_table.h   - Tabela de simbolos
  main.cpp         - Programa principal
assets/            - Programas de teste (.ling)
Makefile           - Automacao de build
\end{lstlisting}

\subsection{Fluxo de compila\c{c}\~ao}

O fluxo de dados entre os m\'odulos segue a seguinte sequ\^encia:

\begin{enumerate}
    \item O \texttt{main.cpp} abre o arquivo fonte (\texttt{.ling}) e configura o \textit{Flex};
    \item O \textit{Flex} (\texttt{lexer.l}) percorre o c\'odigo fonte caractere a caractere, gerando uma lista de tokens;
    \item O \texttt{main.cpp} adiciona um token sentinela \texttt{END\_OF\_FILE} ao final da lista;
    \item O \texttt{Parser} (\texttt{parser.h}) consome a lista de tokens, valida a gram\'atica e preenche a tabela de s\'imbolos;
    \item O resultado \'e exibido: sucesso com tabela de s\'imbolos, ou mensagem de erro.
\end{enumerate}

\section{An\'alise l\'exica}

\subsection{Defini\c{c}\~ao dos tokens}

Os tokens s\~ao definidos no arquivo \texttt{token.h} atrav\'es da enumera\c{c}\~ao \texttt{TokenType}, organizada em 5 categorias:

\begin{itemize}
    \item \textbf{Palavras-chave} (21 tokens): \texttt{class}, \texttt{public}, \texttt{static}, \texttt{void}, \texttt{main}, \texttt{String}, \texttt{return}, \texttt{int}, \texttt{boolean}, \texttt{if}, \texttt{else}, \texttt{while}, \texttt{System}, \texttt{out}, \texttt{println}, \texttt{true}, \texttt{false}, \texttt{this}, \texttt{new}, \texttt{extends}, \texttt{length};

    \item \textbf{Operadores} (8 tokens): \texttt{\&\&}, \texttt{<}, \texttt{>}, \texttt{+}, \texttt{-}, \texttt{*}, \texttt{=}, \texttt{!};

    \item \textbf{Delimitadores} (9 tokens): \texttt{(}, \texttt{)}, \texttt{[}, \texttt{]}, \texttt{\{}, \texttt{\}}, \texttt{;}, \texttt{,}, \texttt{.};

    \item \textbf{Literais e identificadores}: \texttt{ID} (nomes) e \texttt{NUMBER} (inteiros);

    \item \textbf{Controle}: \texttt{END\_OF\_FILE} e \texttt{UNKNOWN}.
\end{itemize}

Cada token \'e representado pela estrutura:

\begin{lstlisting}[style=cpp]
struct Token {
    TokenType type;     // Categoria do token
    std::string lexeme; // Texto original no codigo fonte
    int line;           // Linha onde foi encontrado
};
\end{lstlisting}

\subsection{Implementa\c{c}\~ao com Flex}

O analisador l\'exico \'e implementado no arquivo \texttt{lexer.l} utilizando o gerador \textit{Flex}. O \textit{Flex} gera um aut\^omato finito determin\'istico (AFD) a partir de regras de express\~oes regulares.

\subsubsection{Pr\'e-processamento integrado}

O lexer atua simultaneamente como pr\'e-processador, removendo coment\'arios (de linha e de bloco) e espa\c{c}os em branco.

\subsubsection{Regras de reconhecimento}

As defini\c{c}\~oes de express\~oes regulares reutiliz\'aveis s\~ao:

\begin{lstlisting}[numbers=none]
DIGIT       [0-9]
LETTER      [a-zA-Z]
ID          {LETTER}({LETTER}|{DIGIT}|"_")*
NUMBER      {DIGIT}+
\end{lstlisting}

A prioridade das regras segue a ordem de apari\c{c}\~ao no arquivo: coment\'arios, whitespace, palavras-chave, operadores, delimitadores, identificadores, n\'umeros e erros.

\subsubsection{Tratamento de erros l\'exicos}

Quando o lexer encontra um token inv\'alido, ele calcula a dist\^ancia de Levenshtein entre o token e todas as palavras-chave. Se a dist\^ancia for menor ou igual a 2, sugere a corre\c{c}\~ao mais prov\'avel:

\begin{lstlisting}[numbers=none]
Erro lexico na linha 3: token invalido "clas"
  Sugestao: voce quis dizer "class"?
\end{lstlisting}

\subsection{Sa\'ida da an\'alise l\'exica}

A lista de tokens pode ser visualizada com a flag \texttt{--tokens}:

\begin{lstlisting}[numbers=none]
$ ./build/compiler --tokens assets/prog-factorial.ling
<CLASS, "class", line 1>
<ID, "Factorial", line 1>
<LBRACE, "{", line 1>
<PUBLIC, "public", line 2>
...
\end{lstlisting}

\section{An\'alise sint\'atica}

\subsection{M\'etodo utilizado}

O analisador sint\'atico \'e implementado como um \textit{parser} recursivo descendente no arquivo \texttt{parser.h}. Cada regra da gram\'atica LL(1) corresponde a uma fun\c{c}\~ao no c\'odigo. O parser utiliza \textit{lookahead} de 1 token para decidir qual produ\c{c}\~ao seguir.

\subsection{Gram\'atica da linguagem}

A gram\'atica LL(1) completa utilizada pelo parser foi apresentada na Se\c{c}\~ao 3.2. Cada fun\c{c}\~ao do parser corresponde diretamente a uma produ\c{c}\~ao dessa gram\'atica.

\subsection{Conjuntos FIRST}

Para decidir qual produ\c{c}\~ao seguir, o parser verifica os conjuntos FIRST:

\begin{itemize}
    \item FIRST(Type) = \{ \texttt{int}, \texttt{boolean}, \texttt{ID} \}
    \item FIRST(Cmd) = \{ \texttt{\{}, \texttt{if}, \texttt{while}, \texttt{System}, \texttt{ID} \}
    \item FIRST(Op) = \{ \texttt{\&\&}, \texttt{<}, \texttt{>}, \texttt{+}, \texttt{-}, \texttt{*} \}
    \item FIRST(PrimExp) = \{ \texttt{new}, \texttt{!}, \texttt{(}, \texttt{true}, \texttt{false}, \texttt{ID}, \texttt{NUMBER}, \texttt{this} \}
\end{itemize}

\subsection{Tratamento de erros sint\'aticos}

Ao encontrar um token inesperado, o parser reporta a linha, o token esperado e o encontrado:

\begin{lstlisting}[numbers=none]
Erro sintatico na linha 5: esperado SEMICOLON mas encontrou RBRACE
\end{lstlisting}

\section{Tabela de s\'imbolos}

\subsection{Estrutura}

A tabela de s\'imbolos \'e implementada no arquivo \texttt{symbol\_table.h} como um vetor sequencial. Cada entrada cont\'em:

\begin{itemize}
    \item \texttt{name} --- nome do identificador;
    \item \texttt{type} --- tipo associado;
    \item \texttt{scope} --- escopo de declara\c{c}\~ao;
    \item \texttt{line} --- linha no c\'odigo fonte;
    \item \texttt{category} --- classifica\c{c}\~ao (CLASS, METHOD, INSTANCE\_VAR, LOCAL\_VAR, PARAMETER).
\end{itemize}

\subsection{Preenchimento}

A tabela \'e preenchida durante a an\'alise sint\'atica. Sempre que o parser reconhece uma declara\c{c}\~ao, insere um registro com o escopo apropriado:

\begin{itemize}
    \item Classe $\rightarrow$ escopo \texttt{global}
    \item M\'etodo $\rightarrow$ escopo \'e o nome da classe
    \item Vari\'avel local/par\^ametro $\rightarrow$ escopo \'e \texttt{Classe.metodo}
\end{itemize}

\subsection{Sa\'ida formatada}

Ao final de uma an\'alise bem-sucedida, a tabela \'e impressa:

\begin{lstlisting}[numbers=none]
=== Tabela de Simbolos ===
Nome                Tipo           Escopo              Linha  Categoria
----------------------------------------------------------------------
Factorial           class          global              1      classe
main                void           Factorial           2      metodo
args                String[]       Factorial.main      2      parametro
Fac                 class          global              7      classe
ComputeFac          int            Fac                 9      metodo
num                 int            Fac.ComputeFac      9      parametro
num_aux             int            Fac.ComputeFac      10     var. local
\end{lstlisting}

\section{Compila\c{c}\~ao e execu\c{c}\~ao}

\subsection{Pr\'e-requisitos}

\begin{itemize}
    \item \texttt{g++} (C++17)
    \item \texttt{flex}
    \item \texttt{make}
\end{itemize}

Instala\c{c}\~ao no Ubuntu/WSL:

\begin{lstlisting}[numbers=none]
sudo apt install build-essential flex
\end{lstlisting}

\subsection{Processo de build}

\begin{lstlisting}[numbers=none]
make all        # Compila o projeto
make clean      # Remove artefatos de build
\end{lstlisting}

\subsection{Execu\c{c}\~ao}

\begin{lstlisting}[numbers=none]
./build/compiler programa.ling            # Analise completa
./build/compiler --tokens programa.ling   # Apenas tokens
\end{lstlisting}

\subsection{Sa\'idas poss\'iveis}

\begin{itemize}
    \item \textbf{Sucesso}: mensagem de confirma\c{c}\~ao seguida da tabela de s\'imbolos;
    \item \textbf{Erro l\'exico}: mensagem com linha e sugest\~ao de corre\c{c}\~ao;
    \item \textbf{Erro sint\'atico}: mensagem com linha, token esperado e token encontrado.
\end{itemize}

\section{Considera\c{c}\~oes finais}

O compilador implementado cobre as tr\^es primeiras fases de um compilador tradicional: an\'alise l\'exica, an\'alise sint\'atica e constru\c{c}\~ao da tabela de s\'imbolos. A integra\c{c}\~ao entre o \textit{Flex} e um parser recursivo descendente manual permite controle total sobre as mensagens de erro e a constru\c{c}\~ao incremental da tabela de s\'imbolos.

Como trabalhos futuros, o compilador pode ser estendido com:
\begin{itemize}
    \item An\'alise sem\^antica (verifica\c{c}\~ao de tipos, vari\'aveis n\~ao declaradas);
    \item Gera\c{c}\~ao de c\'odigo intermedi\'ario;
    \item Otimiza\c{c}\~ao e gera\c{c}\~ao de c\'odigo alvo.
\end{itemize}

\end{document}
